% ======================================================================
% EXTENDED DATA / SUPPLEMENTARY — captions and inventory only.
% ======================================================================
\section*{Extended Data and Supplementary Information}

\paragraph{Extended Data Fig. 1 $\vert$ Full geometry and control matrix.}
All six cilia geometries against flat film, short pillars, non-magnetic cilia
and bare sensor; per-condition sensitivity, hysteresis and response time.
\dataneeded{full matrix values}

\paragraph{Extended Data Fig. 2 $\vert$ Decoding validation.}
Train/validation/test split by position, cross-session results and model
comparison. \dataneeded{tables/curves}

\paragraph{Extended Data Fig. 3 $\vert$ Artifact suppression.}
No-contact motion artifacts, reference-sensor compensation and before/after CMR.
\dataneeded{values}

\paragraph{Extended Data Fig. 4 $\vert$ 44-node characterization.}
Per-node baseline, gain, noise, orientation and dead channels.
\dataneeded{node table}

\paragraph{Extended Data Fig. 5 $\vert$ 1$\times$4 strip.}
1D localization, sliding and slip detection. \dataneeded{values}

\paragraph{Extended Data Fig. 6 $\vert$ Full airflow results.}
All wind-tunnel speeds, poses and dynamic profiles. \dataneeded{values}

\paragraph{Extended Data Fig. 7 $\vert$ Full pulse results.}
All subjects, sessions and positions with ECG (and PPG) references.
\dataneeded{values}

\section*{Preliminary single-sensor data (internal; to be superseded)}

\begin{figure}[h]
\centering
\includegraphics[width=\textwidth]{fig5d_pulse_demo.pdf}
\caption*{\textbf{Preliminary Fig.~5d candidate $\vert$ Pulse micro-vibration
demonstration (real data, single MLX90393, L8/d0.8 cilia, 500~Hz).}
Left: raw and band-passed (0.5--10~Hz) $B_z$ at the radial artery.
Middle: beat detection on the band-passed trace. Right: beat-aligned
individual beats (grey) and their mean ($n=95$ beats, peak
$\approx$3.6~$\mu$T), showing a coherent systolic waveform with a secondary
shoulder. Qualitative demonstration only: no synchronized ECG exists for this
recording, so no heart-rate accuracy or timing claim is made
(notes/pulse\_data\_verdict.md).}
\end{figure}

\begin{figure}[h]
\centering
\includegraphics[width=\textwidth]{edx_wind_preliminary.pdf}
\caption*{\textbf{Preliminary airflow data $\vert$ Why the 2025-07-17 session
cannot yield a $\Delta B$--$v$ calibration.} \textbf{(a)}~Band-limited RMS is
non-monotonic in speed. \textbf{(b)}~The response concentrates in narrowband
peaks at fan-rotation harmonics, indistinguishable from fan-motor magnetic
pickup without a bare-sensor control; the broadband floor is
speed-independent. \textbf{(c)}~The steady DC shift follows trial order
(drift and a setup discontinuity at trial~9), not wind speed. Re-measurement
SOP: notes/acquisition\_SOPs.md, SOP-4.}
\end{figure}

\begin{figure}[h]
\centering
\includegraphics[width=0.85\textwidth]{fig2e_geometry_sensitivity.pdf}
\caption*{\textbf{Preliminary Fig.~2e candidate $\vert$ Geometry-dependent
sensitivity (real data, nine $(L,d)$ combinations, manual loading at
1/2/3~mm depths, ATI force reference).} Median plateau sensitivity
$S=|\Delta B|/F$ per geometry under tangential (left) and normal (right)
loading; per-cell values in mT\,N$^{-1}$. The $(8.0,0.8)$ tangential cell and
one $(8,1.0)$-adjacent artifact point are excluded per QC. The largest median
response occurs for the longest, thickest cilia ($L=8$~mm, $d=1.0$~mm).
Manual-loading depth variability (CV 13.4\%) limits absolute calibration; the
robot-controlled campaign (SOP-1) supersedes these values for force claims.}
\end{figure}

\begin{figure}[h]
\centering
\includegraphics[width=0.85\textwidth]{fig2f_dynamics_durability.pdf}
\caption*{\textbf{Preliminary Fig.~2f candidate $\vert$ Durability and
dynamic response (real data).} \textbf{(a)}~Checkpoint sensitivity
$\Delta B_z/\Delta F_z$ over 150{,}000 compression cycles for three
independent samples; open symbols mark checkpoints with $>$9\% corrupted
acquisition packets (low confidence). Reliable checkpoints show no
significant degradation (within $\pm$4\% from 50k to 150k cycles).
\textbf{(b)}~Distribution of 10--90\% response times across 81 tap events
(fastest 25~ms, median 50~ms).}
\end{figure}

\begin{figure}[h]
\centering
\includegraphics[width=0.9\textwidth]{fig3c_event_alignment.pdf}
\caption*{\textbf{Preliminary Fig.~3 material $\vert$ Force--magnetic
correspondence (real data, manual loading, ATI reference).}
\textbf{(a)}~Single load--hold--unload event: compressive force $-F_z$ and
magnetic response $-\Delta B_z$ (sign conventions on axes), peak 370~mN
$\leftrightarrow$ 1009~$\mu$T, Pearson $r=0.95$. \textbf{(b)}~Five
consecutive cycles, $r=0.82$; amplitude variation reflects manual depth
variability.}
\end{figure}

\begin{figure}[h]
\centering
\includegraphics[width=0.9\textwidth]{fig3c_BF_calibration.pdf}
\caption*{\textbf{Preliminary Fig.~3 material $\vert$ $\Delta B$--$F$
relation (real data, manual loading).} \textbf{(a)}~Loading-branch
trajectory of a single press (force relaxation gives the non-monotonic
path). \textbf{(b)}~Plateau amplitudes across 120 manual trials (3 samples
$\times$ 2 target depths): pooled linear fit $k=2.0\times10^3$~$\mu$T/N,
$R^2=0.37$ --- the modest $R^2$ reflects manual depth variability and is
expected to improve substantially under the robot-controlled protocol
(SOP-1).}
\end{figure}

\begin{figure}[h]
\centering
\includegraphics[width=0.9\textwidth]{fig2f_response_event.pdf}
\caption*{\textbf{Preliminary Fig.~2f material $\vert$ Dynamic response and
repeatability (real data).} \textbf{(a)}~Representative median-amplitude tap
event (high-rate acquisition): 10--90\% rise 72~ms, 90--10\% recovery 35~ms
(7-ms moving average; values are smoothing-dependent within $\sim$30--80~ms;
the fastest of 81 events is 25~ms). \textbf{(b)}~Five aligned manual cycles
and mean. \textbf{(c)}~Peak amplitudes across 20 manual trials, CV~13.4\%.}
\end{figure}

\begin{figure}[h]
\centering
\includegraphics[width=0.9\textwidth]{fig2f_cyclic_stability.pdf}
\caption*{\textbf{Preliminary Fig.~2f material $\vert$ Within-checkpoint
cyclic stability of a $5\times5$ cilia array (real data).}
\textbf{(a)}~Baseline-subtracted $B_x$ during machine-cycled loading
(period 0.512~s, 1.95~Hz; 20-s excerpt) at the 10{,}000-cycle checkpoint.
\textbf{(b)}~Per-cycle peak-to-peak amplitude across all 9{,}999 recorded
cycles: median 53.3~$\mu$T, drift between first and last 10\% of cycles
$-2.5$\%.}
\end{figure}

\begin{figure}[h]
\centering
\includegraphics[width=0.62\textwidth]{fig2d_stress_strain_direction.pdf}
\caption*{\textbf{Fig.~2d candidate $\vert$ Stress--strain by magnetization
direction (real data, DIN~EN~ISO~527-1).} Engineering stress (load divided by
each specimen's cross-section $A_0$) versus strain, shown to 100\% strain.
Thin lines: individual specimens; bold lines: group means (averaged only
where all group specimens have data). $E_t$ = linear-fit slope over the
20--40\% strain interval (dimensionless strain), mean\,$\pm$\,s.d.\ from the
machine summary: non-mag.\ $0.72\pm0.11$, $x$ $0.84\pm0.09$, $y$
$0.73\pm0.11$, $z$ $0.68\pm0.11$~MPa. The $x$-magnetized group (aligned with
the tensile axis) shows a trend towards higher stiffness (2 valid
specimens---trend only, no significance claim).}
\end{figure}

\begin{figure}[h]
\centering
\includegraphics[width=0.62\textwidth]{fig2d_stress_strain_wtpct.pdf}
\caption*{\textbf{Fig.~2a/2d candidate $\vert$ Composition series:
Ecoflex~00-30 with 50--80~wt\% NdFeB (real data).} Full curves to specimen
break. 50--70~wt\% remain compliant ($E_t$ $0.130$--$0.152$~MPa) with
elongation above 400\%; 80~wt\% stiffens roughly sixfold
($0.836\pm0.095$~MPa) with elongation reduced to $\approx$150\%,
defining the compliance--magnetic-loading trade-off for cilia design. Note:
the tester export's per-specimen channel is already stress (kPa), verified
against the summary sheet.}
\end{figure}

\begin{figure}[h]
\centering
\includegraphics[width=0.95\textwidth]{fig2_composite.png}
\caption*{\textbf{Figure 2 assembly draft (raster; vector originals ship
separately).} \textbf{(a)}~Composition series (Ecoflex~00-30, 50--80~wt\%
NdFeB). \textbf{(b)}~Fabrication workflow (team schematic; fix ``Vaccum''
$\rightarrow$ Vacuum and ``Demode'' $\rightarrow$ Demould in the vector
redraw). \textbf{(c)}~SEM micrographs with on-image scale bars: c1 array
overview (400~$\mu$m), c2 cilia close-up (300~$\mu$m), c3 single-cilium
cross-section (100~$\mu$m bar; measured width 402.5~$\mu$m for nominal
$d=0.5$~mm). \textbf{(d)}~Stress--strain by magnetization
direction. \textbf{(e)}~Geometry--sensitivity heatmaps. \textbf{(f)}~Dynamic
response, repeatability and 150k-cycle durability (sub-panels f1--f5).}
\end{figure}

\begin{figure}[h]
\centering
\includegraphics[width=\textwidth]{fig3_quad_main.pdf}
\caption*{\textbf{Fig.~3 candidate (main-figure draft) $\vert$ The local
$2\times2$ tri-axis unit separates loading directions.}
\textbf{(a)}~Schematic: sensor layout (top view; cilia patch footprint
inferred from no-load baselines) and machine-driven protocols (shear trials:
1~mm press then $\pm$2~mm shear $\times\sim$20; $z$ trials: 2~mm press
$\times$20; usable trials $x$ 17 of 20, $y$ 19 of 21, $z$ 20 of 20 after
alignment quality control). \textbf{(b)}~Representative $x$-shear trial:
baseline-subtracted 12-channel response (four sensors $\times$
$B_x,B_y,B_z$), left/right shear stages shaded; the four sensors respond
with distinct, spatially differential signatures. \textbf{(c)}~Fingerprint
matrix: trial-averaged stage response per channel for the five loading
conditions (nominal logger $\mu$T). \textbf{(d)}~Per-cycle 12-channel
samples in the LDA plane (visualization projection) and leave-one-trial-out
confusion matrix: 100.0\% (1120/1120 cycles, 56 trials; shuffled-label
control at chance). Two caveats: features are built from protocol stage
labels (shear cycles as right$-$left differentials, $z$ as press-plateau
means), so this is class separability rather than blind real-time decoding;
and each class was acquired as a contiguous same-day block, so
trial-held-out validation cannot fully exclude block-level drift as a cue.
\textbf{(e)}~$z$-press sensitivity of an isolated single sensor (ctrl)
versus the four array nodes; the spread reflects patch--sensor offset
geometry (equal logger gain assumed). Force amplitude is not decodable under
this displacement-controlled protocol ($|r|\leq0.23$; see analysis notes)
--- force decoding awaits the force-controlled xArm campaign.}
\end{figure}

\paragraph{Supplementary Videos.}
(1) $2\times2$ inverse decoding; (2) bubble touch; (3) 44-node real-time map;
(4) no-contact artifact test; (5) wind-tunnel pose change;
(6) ECG-synchronized pulse.
